% Part: sets-functions-relations
% Chapter: relations
% Section: reflections
%
\documentclass[../../../include/open-logic-section]{subfiles}

\begin{document}

\olfileid[pt]{sfr}{rel}{ref}
\olsection{Reflexões Filosóficas}

Na \olref[set]{sec}, definimos relações como certos conjuntos. Devemos
fazer uma pausa e colocar uma questão filosófica rápida: o que é que tal
definição \emph{faz}? É extremamente duvidoso que devamos querer
dizer que \emph{descobrimos} fatos metafísicos de identidade; que, por
exemplo, a relação de ordem em $\Nat$ \emph{acabou por ser} o conjunto
$R=  \Setabs{\tuple{n,m}}{n, m \in \Nat\text{ e }
n < m}$ que definimos na \olref[set]{sec}. Eis três
razões para isso. 

Primeiro: na \olref[set][pai]{wienerkuratowski}, definimos $\tuple{a, b} =
\{\{a\}, \{a, b\}\}$. Consideremos em vez disso a definição $\lVert a,
b\rVert = \{\{b\}, \{a, b\}\} = \tuple{b,a}$. Quando $a \neq b$, temos
que $\tuple{a, b} \neq \lVert a,b\rVert$. Mas poderíamos com a mesma
facilidade ter tomado $\lVert a,b\rVert$ como a nossa definição de par ordenado,
em vez de $\tuple{a,b}$. Ambas as definições teriam funcionado igualmente
bem. Assim, temos dois candidatos igualmente bons a ``ser'' a relação de
ordem nos números naturais, nomeadamente:
\begin{align*}
		R &= \Setabs{\tuple{n,m}}{n, m \in \Nat \text{ e }n < m}\\
		S &= \Setabs{\lVert n,m\rVert}{n, m \in \Nat \text{ e }n < m}.
\end{align*}
Como $R \neq S$, pela extensionalidade, é claro que não podem
\emph{ambos} ser idênticos à relação de ordem em~$\Nat$. Mas seria
meramente arbitrário, e por isso um pouco embaraçoso, afirmar que $R$
em vez de $S$ (ou vice-versa) \emph{é} a relação de ordenação, como
questão de fato. (Isto é um caso muito simples de um argumento contra o
reducionismo em teoria dos conjuntos que Benacerraf tornou famoso em
\citeyear{Benacerraf1965}. Voltaremos a ele várias vezes.)

Segundo: se pensamos que \emph{toda} a relação deve ser identificada
com um conjunto, então a própria relação de pertinência a conjuntos, $\in$, deveria
ser um conjunto particular. De fato, teria de ser o conjunto
$\Setabs{\tuple{x,y}}{x \in y}$. Mas será que este conjunto existe? Dado
o Paradoxo de Russell, é uma afirmação não trivial que tal conjunto exista.
De fato, \oliflabeldef{cumul:::part}{a teoria dos conjuntos que desenvolvemos na
\olref[cumul][][]{part} irá \emph{negar} a existência deste
conjunto.\footnote{Antecipando, eis por que. Por redução ao absurdo, suponhamos que $I =
\Setabs{\tuple{x,y}}{x \in y}$ existe. Então $\bigcup \bigcup I$ é o
conjunto universal, o que contradiz
\olref[sfr][z][sep]{thm:NoUniversalSet}.}}{é possível
desenvolver a teoria dos conjuntos de modo rigoroso como teoria axiomática, e essa
teoria de fato nega a existência deste conjunto.}
Logo, mesmo que algumas relações possam ser tratadas como conjuntos, a relação de
pertinência a conjuntos terá de ser um caso especial.

Terceiro: quando ``identificamos'' relações com conjuntos, dissemos que nos
permitiríamos escrever $Rxy$ para $\tuple{x,y} \in R$. Isto está bem,
desde que a relação de pertinência, ``$\in$'', seja tratada \emph{como}
um predicado. Mas se pensamos que ``$\in$'' denota um certo tipo de
conjunto, então a expressão ``$\tuple{x,y} \in R$'' consiste apenas em
três termos singulares que denotam conjuntos: ``$\tuple{x,y}$'',
``$\in$'', e ``$R$''. E tal lista de nomes não é mais capaz de
exprimir uma proposição do que a sequência sem sentido: ``o copo porta-canetas
a mesa''. De novo, mesmo que algumas relações possam ser tratadas como conjuntos, a
relação de pertinência a conjuntos tem de ser um caso especial. (Isto reúne
num só argumento uma versão simples do paradoxo do conceito \emph{cavalo} de Frege, e
uma objeção célebre que Wittgenstein dirigiu a Russell.)

Então onde nos deixa isto? Bem, não há nada de \emph{errado} em
dizermos que as relações nos números são conjuntos. Só temos de
compreender o espírito em que tal observação é feita. Não enunciamos
um fato metafísico de identidade. Limitamo-nos a notar que, em certos
contextos, podemos (e iremos) \emph{tratar} (certas) relações como
certos conjuntos.

\end{document}
